\scanpage{32}

\mspage{15}

Considérons un module libre \inserted{$M$,} \deleted{\unsure{de r|de}} \deleted{\unsure{\illegible{1}|de}} rang $2$ sur l'anneau $\Lambda=\Lambda_0$ $(\simeq\mathbb{Z}$ ou $\mathbb{Z}/n\mathbb{Z}$, $n\geq3)$, \deleted{et considérons} muni d'une jauge $e_M\in(\det M)^*$, d'où un \unsure{$2$|$1$}-polyèdre régulier triangulé $SP_M$.

Choisissons un repère \deleted{direct} de $SP_M$, \deleted{qui} \deleted{est donc} \inserted{ou} ce qui revient au même, \unsure{une|un} triplet d'éléments distincts de $S_M=M^*/{\pm1}$

\[
(s,t,u)\quad\text{tel que }\deleted{\unsure{\}}}\{s,t,u\}\in SF_M\tag*{(77)}
\]

i.e. tel qu'il existe

\[
\left\{\begin{aligned}
&(x,y,z)\in B_{\deleted{\unsure{\mathrm{sp}}}}(M)\qquad
\left[\begin{aligned}
&\text{donc}\quad\deleted{x+y}\ \inserted{(x,y)\in\operatorname{Rep}(M)}\\
&\qquad x+y+z=0\quad\text{i.e.}\quad z=-(x+y)\\
&\qquad\qquad\deleted{\text{\unsure{d'où}}}
\end{aligned}\right.\\
&\text{avec}\\
&s=\dot{x},\ t=\dot{y},\ u=\dot{z}=(x+y)^{\unsure{\bullet|\cdot}},\quad x\wedge y\in\{\pm e_M\}
\end{aligned}\right.\tag*{(78)}
\]

le repère correspondant étant

\[
r=(s,a,f),\quad\text{avec}\quad a=\{s,t\},\quad f=\{s,t,u\}\;.\tag*{(79)}
\]

\deleted{Le triplet} Supposons que le repère soit direct, i.e la base $(x,y)$ soit directe

\[
x\wedge y=e_M\;.\tag*{(80)}
\]

On va expliciter les automorphismes $\tau_0^r,\tau_1^r,\tau_\infty^r$ du polyèdre régulier $SF_M$, qui correspondent au repère \inserted{(direct)} choisi $r$ (79), i.e. à la base symplectique $(x,y)$ de $M$. (Le choix \deleted{direct} du repère direct revient à celui du choix de la base symplectique, \uline{modulo} signe.)

% archive sheet number: 31
